\section{三价压缩与矩阵乘法指数}\label{sec:shaving}

每次应用引理~\ref{lem:macro} 都会缩小每个副本的规模，同时提高最大元数。我们在规模上的节省仍超过较大查找表的求值费用时停止。此后通过两项操作恢复原域上的三价网络。利用引理~\ref{lem:batch-table} 实施批量求值，再用一次坐标投影消去局部恒等式中使用的扩域。

\begin{lemma}[三价网络的下降]\label{lem:table-descent}
设 $K/k$ 为次数 $D$ 有限的扩张。若规模为 $s$ 的三价 $K$-网络表示一个分量属于 $k$ 的张量，则该张量也可以由一个规模为 $O(sD^2)$ 且外部位置相同的三价 $k$-网络表示。
\end{lemma}

\begin{proof}
选取 $K$ 的一组 $k$-基，使第一个基元素为 $1$。对 $\lambda\in K$，用 $M_\lambda$ 表示乘以 $\lambda$ 的线性映射在该基下的矩阵。将各站点沿一条链排序，每个站点增加两个端口，分别携带传入和传出的基索引。若其原有端口上的一个赋值具有权重 $\lambda$，则扩大后的查找表返回矩阵 $M_\lambda$ 的指定分量。将链的两端均固定为第一个基索引。固定原有连线的任意赋值后，沿链收缩，所得结果为各权重之积的第一坐标。这个坐标映射是 $k$-线性的。因此对原有内部比特求和得到原张量相应分量的第一坐标，而该分量本就属于 $k$。

用 $\lceil\log_2 D\rceil$ 比特编码每个基索引，将未用编码的权重设为零。扩大后的表至多有 $r=3+2\lceil\log_2 D\rceil$ 个比特。每个以 $r$ 个比特为输入的表均有规模为 $O(2^r)$ 的三价实现。在二进制前缀树的每个节点，计算输入是否具有该前缀的指示值。在深度 $j$ 处，将下一个输入比特复制到 $2^j$ 个节点，并检测每个节点的两个子节点，共需 $O(2^j)$ 个门。给每个叶节点配置一个一元权重，指示值为一时取相应的表项，否则取 $1$。恰有一个叶节点提供选定的表项。将各层费用相加，每个原站点的费用为 $O(2^r)=O(D^2)$。
\end{proof}

\begin{proof}[定理~\ref{thm:shaving} 的证明]
固定一个满足 $\omega_k<9/4$ 的有限域 $k$ 和外部比特的容许参数 $\eta>0$。记 $h=\lceil\log_2 n\rceil$。定理~\ref{thm:sort} 中通常宽度的排序给出指数 $1<p<2$。令 $\rho=p-1$。若元数总共由 $h^\beta$ 增至 $h^{b_*}$，则在最终查表时，每个保留下来的关联所需费用为 $O(h^{\rho b_*})$。选取 $0<\delta\le1/10$ 使 $1-\delta>\rho(1+\delta)$ 成立。先取 $A>0$ 充分小，再取 $\beta>0$ 充分小，使 $b_*=\beta+(1+\delta)A$ 满足
\begin{equation}\label{eq:shaving-margins}
 b_*<\tfrac12,\qquad \rho b_*<\tfrac12\eta,\qquad
 \Delta=(1-\delta)A-\rho b_*>0 .
\end{equation}
正的差额 $\Delta$ 即相对于原网络规模的净节省指数。选取充分大的固定整数 $J$，使 $a=A/J$ 在 $\beta\le b\le1/2$ 的整个范围内一致满足引理~\ref{lem:macro} 的要求，并且 $2a<\eta/2$。

这 $J$ 次应用共用一个扩域 $K/k$ 即可。每个宽度尺度 $T$ 介于 $h^\beta$ 的一个常数倍和 $h$ 的某个固定幂之间。局部恒等式要求一个固定的常数域扩张，并要求域的基数至少达到 $T$ 的某个固定幂。选取 $[K:k]$ 充分大，以涵盖这些幂中最大者，同时使所需的各常数域扩张次数均整除它。由于 $k$ 有限，可以做到在每个尺度下均有 $[K:k]=O(\log h)=O(\log T)$。

设 $s_j$ 统计应用 $j$ 次后的站点数及其总元数之和，记 $d_j$ 为累积副本数，并令 $e_j=d_j e$。初始时 $d_0=1$，$s_0=O(m)$，元数至多为 $3\le h^\beta$，而三价 $k$-表具有一个固定的有限字典。在第 $j<J$ 步使用规模额度
\[
 M_j=\max\{s_j,\lceil e_j h^{2a}\rceil\}.
\]
这样，即使外部比特已占当前规模的大部分，构造仍能继续。此时 $M_j\ge n$，$e_j\le M_j h^{-2a}$，且最大元数至多为 $h^{\beta+(1+\delta)aj}$。若本次应用生成 $r_j$ 个副本，则引理~\ref{lem:macro} 保留简短的表标签，给出 $d_{j+1}=r_jd_j$，并得到
\[
 s_{j+1}\le r_j M_j q,\qquad q=h^{-(1-\delta)a}.
\]
它在第~\ref{sec:full} 节中的构造，对所有输入网络一致满足 $\log_2 r_j=O((\log h)^2)$。用归纳法，将原来的界 $s_0,e_0=O(n^{1.05})$ 乘以固定数目的上述复制因子及 $h$ 的固定幂，即得每次调用时均有 $M_j\le n^{1.05+o(1)}\le n^{1.1}$。

若改看每个外部比特对应的规模 $x_j=s_j/e_j$，则复制因子相消。当 $h$ 充分大时，
\[
 x_{j+1}\le q\max\{x_j,2h^{2a}\},\qquad
 x_J\le \max\{q^J x_0,2q h^{2a}\}.
\]
事实上，$q<1$，所以外部比特额度带来的每一项贡献都以最近一次这样的额度为上界。

设 $u$ 为最终的最大元数。应用引理~\ref{lem:batch-table} 生成 $2^u$ 个副本，并将每个副本的规模乘以 $O((1+u)^\rho)$。所得张量为原 $k$-张量的第 $d$ 次张量幂，其中
\[
 d=2^u d_J,\qquad
 \log_2 d\le h^{b_*}+O\bigl(J(\log h)^2\bigr)=o(h).
\]
因此，对于固定的 $k,\eta$，一致地有 $d=n^{o(1)}$。其 $de$ 个外部位置中的每一个，在查表以及应用引理~\ref{lem:table-descent} 的过程中都得以保留。下降的费用为 $O([K:k]^2)=h^{o(1)}$。综合这些界，$k$ 上三价网络的规模除以 $de$ 至多为
\[
 O(m/e)\,h^{-\Delta+o(1)}
       +h^{2a+\rho b_*+o(1)} .
\]
取 $\gamma=\Delta/2$。当 $n$ 充分大时，严格不等式 $\Delta>\gamma$ 与 $2a+\rho b_*<\eta$ 足以吸收各常数和低阶项。再乘以 $de$，即得式~\eqref{eq:shave}。
\end{proof}

对于定理~\ref{thm:many} 的互异性张量，各个独立的外部位置对应彼此分开的各张列表中的条目。选取 $\eta<\sigma$，由命题~\ref{prop:finitefield} 迭代压缩，直至每个外部比特的体积降至 $O(w^\eta)$，低于所需的 $w^\sigma$。这就排除了在任何有限域上成立 $\omega_k<9/4$ 的可能性。引理~\ref{lem:specialize} 将该下界推广到任意域，再结合相匹配的上界，便得定理~\ref{thm:square}。
